% !TeX root = ../../Book3.tex
% 纲要标题：## 附录 D　群作用、大尺度几何与双曲流形
% 纲要位置：几何学计划纲要.md，第 3426 行。

% 纲要内容开始
% 
% > 本附录证明范围与外部输入：规划任务：以第一卷第5、10、16章及本卷第9—10章为实际准备，完整证明有限生成集改变时字度量的拟等距、所选固有测地度量空间上固有且余紧等距作用下的 Švarc–Milnor 引理，并完成整数格、自由群、树及双曲商的代表计算。双曲性、CAT(0) 几何作用与边界分别固定模型和假设；一般增长、群分裂及几何群论分类不承担完整证明。三维双曲流形只构造明确的完整商与尖点模型，拓扑分解和几何化的准确陈述与第二卷附录A第4节分工；第27.5节不反向调用此处尚未写成的内容。Mostow–Prasad 刚性限曲率归一化为负一、维数至少三、连通完备有限体积的所选流形版本，准确外引并完整推出所选体积不变量结论；不混同二维 Teichmüller 形变与无限体积情形。
% 
% * **字度量与群作用**
%   * 有限生成群的 Cayley 图
%   * 生成集变化与拟等距
%   * 等距作用、轨道映射与基本区域
%   * 固有性、余紧性与 Švarc–Milnor 引理
%   * 图、整数格与双曲商的完整算例
% 
% * **增长与端**
%   * 群的增长函数及生成集无关性
%   * 多项式增长与指数增长的模型
%   * 图与群的端
%   * 拟等距对增长和端的保持
%   * 虚循环群与树作用的比较
%   * 一般增长和群分裂定理的外部输入
% 
% * **双曲性与非正曲率群作用**
%   * 薄三角形与 Gromov 双曲性
%   * 树与双曲空间的验证
%   * 双曲群的边界与群作用
%   * CAT(0) 空间上的几何作用
%   * 平坦与双曲性的条件差异
%   * 边界模型与一般结构定理的外部输入
% 
% * **三维双曲商与边界**
%   * 双曲三空间及理想边界
%   * 离散等距群与 Kleinian 群
%   * 无挠商、完备性与流形条件
%   * 抛物子群、尖点与有限体积
%   * 共形边界、极限集与几何有限性的区别
%   * 显式基本区域与三维商模型
% 
% * **有限体积双曲流形的刚性**
%   * 带度量的流形与基本群数据
%   * Mostow–Prasad 刚性的准确陈述
%   * 维数、完备性与有限体积的作用
%   * 二维模参数与无限体积的边界
%   * 体积不变量及典型推论
%   * 一般证明的外部工具与阅读条件
% 
% ---
% 
% 纲要内容结束
% 正文状态：未撰写；本文件目前是框架。
% 规划任务：以第一卷第5、10、16章及本卷第9—10章为实际准备，完整证明有限生成集改变时字度量的拟等距、所选固有测地度量空间上固有且余紧等距作用下的 Švarc–Milnor 引理，并完成整数格、自由群、树及双曲商的代表计算。双曲性、CAT(0) 几何作用与边界分别固定模型和假设；一般增长、群分裂及几何群论分类不承担完整证明。三维双曲流形只构造明确的完整商与尖点模型，拓扑分解和几何化的准确陈述与第二卷附录A第4节分工；第27.5节不反向调用此处尚未写成的内容。Mostow–Prasad 刚性限曲率归一化为负一、维数至少三、连通完备有限体积的所选流形版本，准确外引并完整推出所选体积不变量结论；不混同二维 Teichmüller 形变与无限体积情形。
% 后续写作分别交代摘要、目标、所需基础、代表算例、习题和小结，不用标题或占位文字宣称已完成。

\chapter{群作用、大尺度几何与双曲流形}
\label{g3:app:D}
\BookChapterNumber{g3:app:D}

待撰写

\section{字度量与群作用}
\label{g3:sec:D01}

待撰写

\subsection{有限生成群的 Cayley 图}
\label{g3:subsec:D01:01}
待撰写

\subsection{生成集变化与拟等距}
\label{g3:subsec:D01:02}
待撰写

\subsection{等距作用、轨道映射与基本区域}
\label{g3:subsec:D01:03}
待撰写

\subsection{固有性、余紧性与 Švarc–Milnor 引理}
\label{g3:subsec:D01:04}
待撰写

\subsection{图、整数格与双曲商的完整算例}
\label{g3:subsec:D01:05}
待撰写

\section{增长与端}
\label{g3:sec:D02}

待撰写

\subsection{群的增长函数及生成集无关性}
\label{g3:subsec:D02:01}
待撰写

\subsection{多项式增长与指数增长的模型}
\label{g3:subsec:D02:02}
待撰写

\subsection{图与群的端}
\label{g3:subsec:D02:03}
待撰写

\subsection{拟等距对增长和端的保持}
\label{g3:subsec:D02:04}
待撰写

\subsection{虚循环群与树作用的比较}
\label{g3:subsec:D02:05}
待撰写

\subsection{一般增长和群分裂定理的外部输入}
\label{g3:subsec:D02:06}
待撰写

\section{双曲性与非正曲率群作用}
\label{g3:sec:D03}

待撰写

\subsection{薄三角形与 Gromov 双曲性}
\label{g3:subsec:D03:01}
待撰写

\subsection{树与双曲空间的验证}
\label{g3:subsec:D03:02}
待撰写

\subsection{双曲群的边界与群作用}
\label{g3:subsec:D03:03}
待撰写

\subsection{CAT(0) 空间上的几何作用}
\label{g3:subsec:D03:04}
待撰写

\subsection{平坦与双曲性的条件差异}
\label{g3:subsec:D03:05}
待撰写

\subsection{边界模型与一般结构定理的外部输入}
\label{g3:subsec:D03:06}
待撰写

\section{三维双曲商与边界}
\label{g3:sec:D04}

待撰写

\subsection{双曲三空间及理想边界}
\label{g3:subsec:D04:01}
待撰写

\subsection{离散等距群与 Kleinian 群}
\label{g3:subsec:D04:02}
待撰写

\subsection{无挠商、完备性与流形条件}
\label{g3:subsec:D04:03}
待撰写

\subsection{抛物子群、尖点与有限体积}
\label{g3:subsec:D04:04}
待撰写

\subsection{共形边界、极限集与几何有限性的区别}
\label{g3:subsec:D04:05}
待撰写

\subsection{显式基本区域与三维商模型}
\label{g3:subsec:D04:06}
待撰写

\section{有限体积双曲流形的刚性}
\label{g3:sec:D05}

待撰写

\subsection{带度量的流形与基本群数据}
\label{g3:subsec:D05:01}
待撰写

\subsection{Mostow–Prasad 刚性的准确陈述}
\label{g3:subsec:D05:02}
待撰写

\subsection{维数、完备性与有限体积的作用}
\label{g3:subsec:D05:03}
待撰写

\subsection{二维模参数与无限体积的边界}
\label{g3:subsec:D05:04}
待撰写

\subsection{体积不变量及典型推论}
\label{g3:subsec:D05:05}
待撰写

\subsection{一般证明的外部工具与阅读条件}
\label{g3:subsec:D05:06}
待撰写
